\documentclass[10pt]{article}
\usepackage[a4paper,margin=1.9cm]{geometry}
\usepackage{amsmath,amssymb}
\usepackage{graphicx}
\graphicspath{{../}{./}}
\usepackage{float}
\usepackage[dvipsnames]{xcolor}
\usepackage{enumitem}
\usepackage{titlesec}
\usepackage[hidelinks]{hyperref}

% --- compact layout ---
\setlist[itemize]{leftmargin=1.35em,itemsep=1pt,topsep=2pt,parsep=0pt}
\setlength{\parindent}{0pt}
\setlength{\parskip}{3pt}
\titlespacing*{\section}{0pt}{6pt}{3pt}
\titlespacing*{\subsection}{0pt}{4pt}{2pt}
\titleformat{\section}{\normalfont\large\bfseries\color{RoyalBlue}}{\thesection.}{0.5em}{}
\titleformat{\subsection}{\normalfont\bfseries}{}{0em}{}

\newcommand{\code}[1]{\texttt{\small #1}}
\newcommand{\file}[1]{\textcolor{RoyalBlue}{\texttt{\small #1}}}

% \docheader{title}{subtitle}
\newcommand{\docheader}[2]{%
  \begin{center}
    {\LARGE\bfseries #1}\\[2pt]
    {\small #2}
  \end{center}
  \vspace{-4pt}\hrule\vspace{6pt}
}

\begin{document}

\docheader{Adora PINN --- Step 7: Prototype inversion}%
{As-built notes \quad$\cdot$\quad Ivan Milic \quad$\cdot$\quad 2026-09-20 \quad$\cdot$\quad step 7 of pinn/plan.md}

Step 7 is an end-to-end inversion of the observed testset (step 5): starting from the pretrained FAL-C field
(warm start), optimise the field weights to fit the Stokes observations, regularised by hydrostatic
equilibrium and the top-boundary pressure. Files \file{pinn/invert.py}, \file{pinn/step7\_inversion.ipynb};
result saved to \file{/dat/milic/adora\_pinn\_develop/inverted\_field.npz}.

\subsection*{1. Machinery}
\code{invert.run\_inversion} runs optax Adam on the field weights with the total loss
$w_\mathrm{fit}\,\chi^2 + w_\mathrm{he}\,\mathcal L_\mathrm{HE} + w_\mathrm{bnd}\,\mathcal L_\mathrm{bnd}$.
\textbf{Fresh random samples every step} --- fit pixels (\code{m\_fit}), HE collocation points (\code{k\_he})
and boundary pixels (\code{m\_bnd}) --- so the fit is stochastic over the map and HE is enforced over the
whole domain in expectation; a separate FIXED monitor set gives clean loss curves. The forward model is
\code{rt\_forward.synth\_field} (field $\to$ EOS $\to$ piecewise-linear RT), differentiated end-to-end.

\subsection*{2. Continuable training \& scheduler}
\code{run\_inversion} takes either a params pytree (fresh) or a \textbf{checkpoint} dict and returns
\code{\{params, state, key, step, hist, monitor\}}; feeding it back continues training --- the Adam moments,
RNG, LR-schedule step count and accumulated history all carry over (the notebook has a re-runnable
``continue'' cell). The learning rate follows an \textbf{exponential-decay schedule}
(\code{lr}$\cdot$\code{lr\_decay}$^{\,\mathrm{step}/\mathrm{lr\_transition}}$, floored) --- the standard,
open-ended choice for PINNs. Per-Stokes fit weights are \textbf{hand-set} (\code{sigma\_stokes}, default
$(10^{-2},5{\cdot}10^{-3},5{\cdot}10^{-3},5{\cdot}10^{-3})\times\langle I_c\rangle$ --- $Q,U,V$ a few$\times$
tighter than $I$, not orders apart).

\subsection*{3. Results (2000 steps, one H100)}
Misfit RMS$(\text{model}-\text{obs})/\langle I_c\rangle$, initial $\to$ final:
$I\;0.198\!\to\!0.050$, $Q\;1.2{\cdot}10^{-3}\!\to\!1.2{\cdot}10^{-3}$, $U\;1.2{\cdot}10^{-3}$,
$V\;9.3{\cdot}10^{-3}\!\to\!8.3{\cdot}10^{-3}$. The notebook shows the loss/LR curves, obs-vs-inverted
$I,Q,U,V$ profiles, and the recovered continuum, Doppler $v_z$ and longitudinal $B_z$ maps. Stokes $I$ (the
thermodynamics + velocity) is recovered well; the polarisation is only partly fit.

\subsection*{4. Known limitation \& next}
The recovered maps show \textbf{stripes} and $V/Q/U$ recovery is limited --- diagnosed separately in
\file{adora\_pinn\_7\_debug.pdf}: the random Fourier modes show through wherever the data under-constrains the
field, and this is best cured with a \textbf{spatial-smoothness regularizer} (plus moderate $\sigma$, denser
sampling, and regularizing/freezing weak channels). That regularizer is the natural next addition to
\file{losses.py}. Files uncommitted.

\end{document}
